% SEGMENT OLP-0577-S01
% Part: set-theory
% Chapter: ord-arithmetic
% Section: using-addition

\documentclass[../../../include/open-logic-section]{subfiles}

\begin{document}

\olfileid{sth}{ord-arithmetic}{using-addition}
\olsection{வரிசையெண் கூட்டலைப் பயன்படுத்துதல்}

வரிசையெண் கூட்டலைப் பயன்படுத்தி, \olref[spine][rank]{defnsetrank}
இல் உள்ள பொருளில் பல்வேறு கணங்களின் தரநிலைகளை வெளிப்படையாகக்
கணக்கிடலாம்:

\begin{lem}\ollabel{rankcomputation}
$\setrank{A} = \alpha$ மற்றும் $\setrank{B} = \beta$ எனில்:
\begin{enumerate}
  \item\ollabel{exrankpow} $\setrank{\Pow{A}} = \alpha\ordplus 1$
  \item\ollabel{exrankpair} $\setrank{\{A, B\}} = \max(\alpha,
  \beta) \ordplus 1$
  \item\ollabel{exrankcup} $\setrank{A \cup B} = \max(\alpha,
  \beta)$
  \item\ollabel{exranktuple} $\setrank{\tuple{A,B}} = \max(\alpha,
  \beta) \ordplus  2$
  \item\ollabel{exranktimes} $\setrank{A \times B} \leq
  \max(\alpha, \beta) \ordplus  2$
  \item\ollabel{exrankunion} $\alpha$ வெறுமையாகவோ எல்லை
  வரிசையெண்ணாகவோ இருந்தால் $\setrank{\bigcup A} = \alpha$;
  $\alpha = \gamma\ordplus 1$ எனில் $\setrank{\bigcup A} = \gamma$
\end{enumerate}
\end{lem}

\begin{proof}
நிறுவல் முழுவதும் \olref[spine][rank]{ranksupstrict} ஐ மீண்டும்
மீண்டும் பயன்படுத்துகிறோம்.

\emph{\olref{exrankpow}.} $x \subseteq A$ எனில்
$\setrank{x} \leq \setrank{A}$. ஆகவே
$\setrank{\Pow{A}} \leq \alpha \ordplus  1$.
குறிப்பாக $A \in \Pow{A}$ என்பதால்
$\setrank{\Pow{A}} = \alpha \ordplus  1$.

\emph{\olref{exrankpair}.} \olref[spine][rank]{ranksupstrict} இன் படி.

\emph{\olref{exrankcup}.} \olref[spine][rank]{ranksupstrict} இன் படி.

\emph{\olref{exranktuple}.} \olref{exrankpair} ஐ இருமுறை
பயன்படுத்தினால்.

\emph{\olref{exranktimes}.}
$A \times B \subseteq \Pow{\Pow{A \cup B}}$ என்பதைக் கவனித்து,
\olref{exrankpow} ஐ இருமுறை பயன்படுத்துக.
% SOURCE CORRECTION TA-STH-025: two powersets require exrankpow twice, not exranktuple.

\emph{\olref{exrankunion}.} $\alpha = \gamma\ordplus 1$ எனில்,
$\setrank{c} = \gamma$ ஆக அமையும் ஏதோ ஒரு $c \in A$ உண்டு;
$A$ இன் எந்த !!{element} உம் இதைவிட உயர்ந்த தரநிலை
கொண்டிருக்காது. எனவே $\setrank{\bigcup A} = \gamma$.
$\alpha$ ஓர் எல்லை வரிசையெண் எனில், $\alpha$ வை விடக்
கண்டிப்பாகக் குறைவாக, ஆனால் அதனை எவ்வளவு வேண்டுமானாலும்
நெருங்கும் தரநிலையுடைய !!{element}s $A$ இல் உள்ளன. ஆகவே
$\bigcup A$ இலும் $\alpha$ வை விடக் கண்டிப்பாகக் குறைவாக,
ஆனால் அதனை எவ்வளவு வேண்டுமானாலும் நெருங்கும் தரநிலையுடைய
!!{element}s உள்ளன; அதனால் $\setrank{\bigcup A} = \alpha$.
\end{proof}
\noindent
\olref{exranktimes} இல் ஏன் \emph{சமனின்மை} வருகிறது என்பதைக்
காட்டுவதைப் பயிற்சியாக விடுகிறோம்.

\begin{prob}
$\setrank{A \times B} =
\max(\setrank{A}, \setrank{B})$ ஆக அமையும் $A$, $B$ கணங்களைத்
தருக. $\setrank{A \times B} =
\max(\setrank{A}, \setrank{B}) \ordplus  2$ ஆக அமையும் $A$,
$B$ கணங்களையும் தருக. வேறு தரநிலைகள் சாத்தியமா?
% SOURCE CORRECTION TA-STH-026: restore the missing equality in the second request.
\end{prob}

% SEGMENT OLP-0577-S02
ஒரு வரிசையெண் ``முடிவுறு'' அல்லது ``முடிவுறாதது'' என்பதற்குப்
பொருள் என்ன என்பதற்கான பல நியாயமான கருத்துகள் ஒன்றுக்கொன்று
சமானமானவை என்பதையும் இப்போது காட்டலாம்:

\begin{lem}\ollabel{ordinfinitycharacter}
எந்த வரிசையெண் $\alpha$ க்கும், பின்வருவன சமானமானவை:
\begin{enumerate}
  \item\ollabel{ord:notinomega} $\alpha\notin \omega$;
  அதாவது, $\alpha$ ஓர் இயல் எண் அல்ல
  \item\ollabel{ord:omegaplus} $\omega \leq \alpha$
  \item\ollabel{ord:oneplus} $1 \ordplus  \alpha = \alpha$
  %\alpha \approx \alpha \ordplus 1$
  \item\ollabel{ord:plusone} $\alpha \approx \alpha\ordplus 1$;
  அதாவது, $\alpha$ வும் $\alpha\ordplus 1$ உம் எண்ணொத்தவை
  \item\ollabel{ord:infinite} $\alpha$ டெடெகிண்ட்-முடிவுறாதது
\end{enumerate}
\end{lem}
\noindent

% SEGMENT OLP-0577-S03
ஆக, ஒரு வரிசையெண் முடிவுறுகிறதா முடிவுறாததா என்பதைப் புரிந்துகொள்ள
ஒன்றுக்கொன்று சமானமான ஐந்து வழிகளை நிறுவியுள்ளோம்.

\begin{proof}
\emph{\olref{ord:notinomega} $\Rightarrow$ \olref{ord:omegaplus}.}
முக்கூற்றுப் பண்பின்படி.

\emph{\olref{ord:omegaplus} $\Rightarrow$ \olref{ord:oneplus}.}
$\alpha \geq \omega$ என்க. முடிவிலிக்கடந்த தொகுத்தறிதலின்படி,
$\alpha = \beta \ordplus \gamma$ என்றும் $\beta$ ஓர் எல்லை
வரிசையெண் என்றும் ஆகுமாறு ஏதோ ஒரு மீச்சிறு வரிசையெண் $\gamma$
உண்டு; அது $0$ ஆகவும் இருக்கலாம். இப்போது:
\[
  1 \ordplus \alpha =
  1 \ordplus (\beta \ordplus \gamma) =
  (1 \ordplus \beta) \ordplus \gamma =
  \supstrict_{\delta < \beta} (1 \ordplus  \delta) \ordplus  \gamma =
  \beta \ordplus  \gamma =
  \alpha.
\]
\emph{\olref{ord:oneplus} $\Rightarrow$ \olref{ord:plusone}.}
$f \colon (\alpha \disjointsum 1) \to (1 \disjointsum \alpha)$
என்ற !!a{bijection} இருப்பது தெளிவு.
$1 \ordplus \alpha = \alpha$ எனில்,
$g \colon (1 \disjointsum \alpha) \to \alpha$ என்ற
வரிசை ஓரினச்சார்பு உண்டு. இப்போது $\comp{f}{g}$ ஐக் கருதுக.

\emph{\olref{ord:plusone} $\Rightarrow$ \olref{ord:infinite}.}
$\alpha \approx \alpha \ordplus 1$ எனில்,
$f \colon (\alpha \disjointsum 1) \to \alpha$ என்ற
!!a{bijection} உண்டு. ஒவ்வொரு $\gamma < \alpha$ க்கும்
$g(\gamma) = f(\gamma, 0)$ என வரையறுக்க. இந்த
!!{injection} ஆனது $\alpha$ டெடெகிண்ட்-முடிவுறாதது என்பதைக்
காட்டுகிறது; ஏனெனில்
$f(0,1) \in \alpha \setminus \ran{g}$.

\emph{\olref{ord:infinite} $\Rightarrow$ \olref{ord:notinomega}.}
இது \olref[z][infinity-again]{naturalnumbersarentinfinite} ஆகும்.
\end{proof}

\end{document}
